\section{Complexity ledger}
\label{app:complexity}

Let a dense local matrix have \(O(k)\) rows and columns.  Conventional Gaussian
elimination uses \(O(k^3)\) arithmetic operations and \(O(k^2)\) stored field
elements.  Matrix multiplication and rank reduction have the same classical
cubic upper bound.  Over \(N\) sample bags this yields
\(O(Nk^3)\) operations and, when streamed, \(O(k^2)\) working field elements.

For \(k=P+K\) and \(N,P,K\le n\), monotonicity gives
\begin{align}
 N(P+K)^3
 &\le n(n+n)^3\\
 &=n(2n)^3\\
 &=8n^4.
\end{align}
This inequality has been machine-checked in the Lean development.

\subsection{Conditional rational-matrix bit bound}

Suppose an \(r\times c\) matrix with rational entries of bit length at most
\(L\) is explicitly given.  Clearing denominators produces an integer matrix
with polynomially related bit length.  Fraction-free elimination represents
intermediate entries through minors; determinant bounds then give polynomial
bit lengths in \(r,c,L\).  This supports a polynomial Turing bound for the
linear-algebra phase~\cite{Bareiss1968}.

The premises are not automatically satisfied by symbolic CMP:
\begin{itemize}
  \item a satisfying point is not given;
  \item its coordinates need not be rational;
  \item symbolic entries lie in a polynomial or quotient ring; and
  \item the zeroth-order message may contain several semialgebraic components.
\end{itemize}

\subsection{Checklist for a future bit theorem}

A complete proof of \cref{eq:desired-bit-bound} must state:
\begin{enumerate}
  \item a canonical binary encoding for every message;
  \item polynomial bounds on message length after every merge and projection;
  \item exact algorithms for equality, sign, and emptiness in that encoding;
  \item coefficient-height bounds under every operation;
  \item the number of operations along the tree; and
  \item the cost of the root acceptance test.
\end{enumerate}
Without these items, an \(O(n^4)\) count of abstract ring or field operations is
not a proof of membership in \(P\).
